% ======================================================================
\section{Discussion}\label{sec:discussion}

This section interprets the results. It relates them to the expectations of \cref{sec:framework} and to previous research, discusses what the falling carbon requirement and the decoupling results do and do not show, and draws out implications for countries in each configuration. It ends with the study's limitations and directions for future research.

\subsection{The findings against the expectations}

All four expectations set out in \cref{sec:framework} are borne out, in some cases more starkly than anticipated.

\textbf{E1, few or no countries in the safe and just space,} holds in its strongest form. Under the main specification no country with consumption-based data was in the safe and just space in any year, and every country that crossed the social foundation did so above the boundary. In the alternative specifications, at most \nRobSafeMax{} country appears there in 2011--2022.

\textbf{E2, a falling carbon requirement,} also holds. The requirement fell by a factor of \nReqFactor{} over three decades. But it remained \nReqEndMultiple{} times the 2\textdegree C boundary, and the answer to the open question, whether it has fallen to the boundary, is no.

\textbf{E3, widespread but insufficient decoupling,} holds. The number of absolute decouplers doubled between 2001--2010 and 2011--2022, from \nDecAbsEarly{} to \nDecAbsLate. Only \nDecAboveOnTrack{} of the \nDecAbsAbove{} decouplers above the boundary were on track, and the gap between achieved and required rates is large enough that even optimistic scenarios close it for only a minority.

\textbf{E4, persistent configurations with movement towards the foundation,} holds. Between \nPersistOne\% and \nPersistTwo\% of countries stayed in the same configuration from one period to the next. Almost all movement was towards higher social attainment and, with it, higher emissions.

\subsection{Relation to previous research}

The finding that no country in the sample meets this social foundation within its share of the climate boundary confirms, with more recent data and consumption-based emissions, the conclusion of \citet{oneill2018} and \citet{fanning2022}. Those studies found that no country met basic needs at a globally sustainable level of resource use. This paper uses a narrower set of social indicators and only the climate boundary, so its foundation is easier to meet than theirs. That it still finds the safe and just space empty strengthens the conclusion. The finding that every crossing of the social foundation happened above the boundary adds a dynamic dimension. It is not only that no country is in the safe and just space at a given time. Historically, the path into the social foundation has not passed through it.

The falling carbon requirement is consistent with the long-run decline in the energy and carbon associated with given levels of human development documented by \citet{steinberger2010}, and with the observation of \citet{lamb2015} that the relationship between development and emissions has shifted over time. It is also consistent with the shape of the relationship described by \citet{porcher2026}: social outcomes rise steeply with emissions at low levels and then flatten. The much lower requirement for the health thresholds than for the full foundation echoes the contrast drawn by \citet{steinberger2020} between the loose coupling of life expectancy with emissions and the tight coupling of income with emissions.

The decoupling results extend the conclusion of \citet{vogel2023} beyond high-income countries. They found that 11 high-income countries had absolutely decoupled consumption-based emissions from GDP in 2013--2019, at rates far below those compatible with the Paris Agreement. Here, with per-person emissions and a longer window, \nDecAbsLate{} countries decoupled absolutely in 2011--2022, including many middle-income countries. But the conclusion about speed is unchanged. The median absolute decoupler above the boundary cut emissions by \nDecMedianAchieved\% a year, less than a third of the rate required to reach the boundary by 2050. This is in line with the broader evidence reviewed by \citet{haberl2020}.

\subsection{Interpreting the falling carbon requirement}

The carbon requirement summarises how the probability of meeting the social foundation varies with emissions across countries. Its decline means that, at any given level of emissions, a country in \nLastYear{} was much more likely to meet the foundation than a country at the same emission level in the early 1990s. This is a statement about the cross-country distribution. It has three important qualifications.

First, the requirement is not a causal estimate of the emissions needed to meet basic needs. Countries at the same emission level differ in many ways, including income, institutions, health systems and geography, that also affect social outcomes. The decline in the requirement could reflect changes in any of these, not only in the efficiency with which emissions are converted into social outcomes.

Second, this study does not test why the requirement fell. Plausible contributors include cleaner energy systems, so that a given level of energy services involves fewer emissions; cheaper and more widely available public health technologies, such as vaccines and oral rehydration; and the expansion of electricity access through grids that increasingly include renewables. Testing these explanations would require data on energy systems and health provision that are beyond the scope of this paper.

Third, the level of the requirement depends on the definition of the social foundation. Under the main definition it was \nReqEnd~t in \nLastYear. Without the employment threshold it was \nRobReqNoEmp~t, and with secondary enrolment added it was \nRobReqEdu~t. The employment threshold matters because unemployment rates of above 6\% are common in some upper-middle- and high-income countries, so requiring 94\% employment excludes some countries with high social outcomes and high emissions. Which definition is preferable is partly a normative question. The robust findings are that the requirement has fallen in every specification and remains above the boundary in every specification.

As an illustration only, if the requirement kept falling at the constant proportional rate estimated by the pooled model, it would reach the 2\textdegree C boundary around \nReqReachYear{}. That is well after the 2050 horizon implied by the carbon budget, and the illustration says nothing about the stricter 1.5\textdegree C boundary. There is no reason to expect the rate to stay constant. It could accelerate as clean energy becomes cheaper, or slow as the remaining social shortfalls, concentrated in the poorest countries, prove harder to close.

\subsection{Decoupling: direction versus speed}

The decoupling results show a clear change in direction. In 2011--2022, falling consumption-based emissions with growing incomes became common, rather than exceptional. Falling emissions were generally accompanied by rising life expectancy, so decoupling did not come at the expense of health. These are significant achievements.

The boundary, however, sets a destination, and what matters is the speed of travel. The comparison of achieved and required rates shows a gap that cannot be closed by modest acceleration. Doubling every country's rate would bring only \nScenDouble{} of the \nDecAbsAbove{} to the boundary by 2050. Even if all of them matched the fastest decline any of them achieved over a decade, only \nScenFastest{} would arrive in time. These scenarios are arithmetic illustrations, not forecasts. Future emissions will depend on policy, technology and economic change. But they make clear that reaching the boundary by mid-century would require sustained rates of decline that no country in the sample has yet achieved for a decade.

\subsection{Heterogeneity and the frontier}

The averages hide a wide range. A handful of countries met the social foundation with consumption-based emissions of about two tonnes per person, close to the boundary, while others met it with up to \nMetMaxCons~t. A group of \nNearN{} countries are within 10\% of every threshold without meeting all of them. These differences suggest that the carbon cost of meeting basic needs is not fixed by technology alone, and that the way societies provide for basic needs matters \citep{vogel2021}. They also suggest a practical agenda: understanding what allows the frontier countries to achieve high social outcomes with relatively low emissions, and whether it can be replicated elsewhere. This paper identifies these countries but does not explain their position.

\subsection{Implications for countries in each configuration}

The four configurations imply different tasks.

\paragraph{Foundation met, overshoot.} For these countries, the main task is to reduce emissions while maintaining social outcomes. The decoupling results show that many are reducing emissions, but the median required rate is \nDecSpeedRatio{} times the median achieved rate. The frontier countries, and the fact that life expectancy kept rising in most decoupling countries, suggest that deep reductions need not compromise basic social outcomes. That conclusion does not extend automatically to other dimensions of well-being not measured here.

\paragraph{Shortfall and overshoot.} This is the largest group. These countries already exceed their share of the carbon budget but have not yet met the social foundation. They face both tasks at once. The falling carbon requirement is the most encouraging finding for them: meeting the foundation has become much less carbon-intensive than it was, and the health dimensions in particular are now associated with emissions not far above the boundary.

\paragraph{Shortfall within the boundary.} These countries, mostly in Africa and Asia (\nWithinLastAfrica{} and \nWithinLastAsia{} of the \nWithinLastTotal{} in 2011--2022), have low emissions because consumption is low. Their task is social progress, and the question is how much additional carbon it will require. History suggests that countries leaving this configuration have moved into overshoot. The falling requirement and the frontier countries suggest that the carbon cost of this progress can be much lower than it was for countries that developed earlier, particularly if electricity access is expanded through low-carbon sources.

\paragraph{The safe and just space.} No country in the sample is in it on a consumption basis when the thresholds are applied strictly. Allowing a 5\% shortfall on each indicator brings a few low-emitting Central American countries into it (\cref{sec:res-robust}). The challenge is not to protect a group of successful countries, but to create a path into the space that no country has yet followed.

\subsection{Limitations}

The results should be read in light of the following limitations.

\begin{enumerate}
\item \textbf{One planetary boundary.} The analysis covers only the climate boundary. It says nothing about biosphere integrity, land-system change, freshwater use or biogeochemical flows, where \citet{fanning2022} found similar or larger overshoot. A country within the climate boundary may overshoot others. Adding a material-footprint dimension \citep{oneill2018} would be a natural extension.

\item \textbf{A narrow social foundation.} The foundation has four indicators. Education, income poverty, nutrition, sanitation, equality and democratic quality are not included, because their data are incomplete for the countries and years studied. Meeting this foundation therefore shows that four basic needs are met, not that all basic needs are met, and the four indicators do not capture many dimensions of human well-being. A broader foundation would raise the carbon requirement, as the robustness check with secondary enrolment shows.

\item \textbf{Strict thresholds.} A country meets the foundation only if all four thresholds are met, and exceeds the boundary if its emissions are above \nBoundary~t by any amount. The counts of countries in each configuration therefore depend on the exact threshold values. With a 5\% tolerance on each social threshold, \nRobTolSafeYearN{} countries are in the safe and just space in \nLastYear{} (\cref{sec:res-robust}).

\item \textbf{Sensitivity to the employment threshold.} The level of the carbon requirement depends strongly on whether employment is included. The definition used follows the cited studies, but other reasonable definitions give different levels.

\item \textbf{Coverage of consumption-based emissions.} Consumption-based estimates are not available for \nMissOutN{} of the \nNCountries{} countries. These countries are smaller, poorer and lower-emitting than the main sample (\cref{tab:missing}), so the main sample under-represents countries with social shortfalls within the boundary. The production-based analysis, which covers all countries, gives the same conclusions.

\item \textbf{Uncertainty in consumption-based emissions.} Consumption-based emissions are estimated from multi-regional input--output models and are less certain than territorial emissions, particularly for smaller economies with complex trade structures. The confidence intervals reported here reflect sampling uncertainty across countries, not measurement uncertainty in the emissions estimates.

\item \textbf{The equal-per-capita boundary.} The boundary is an equal share per person of a global budget. Allocations that account for historical emissions or capacity to act would make it stricter for high-income countries \citep{hickel2022,haeyhae2016}, strengthening the conclusions for those countries. The boundary also compares fossil and industrial \COtwo{} with a budget that includes land-use emissions, which slightly understates overshoot for countries with net land-use emissions.

\item \textbf{Associations, not causes.} The carbon requirement and the fixed-effects and long-difference estimates describe associations. They do not identify how much social progress additional emissions cause, or what explains the changes over time.

\item \textbf{Extrapolations are illustrations.} The years to reach the boundary, the year in which the requirement would reach the boundary, and the decoupling scenarios all assume constant rates of change. They illustrate the size of the gaps and are not forecasts.

\item \textbf{Data vintage.} The social indicators come from an extract of the World Development Indicators compiled in early 2026. National statistics are revised regularly. A comparison with the October 2026 release for the countries closest to the thresholds found identical values for life expectancy and unemployment, one small difference for electricity access, and revisions of up to 2 deaths per 1{,}000 in under-5 mortality, none of which changes a classification (\cref{app:repro}). Larger future revisions could change individual values, particularly for recent years.
\end{enumerate}

\subsection{Directions for future research}

Four extensions would strengthen and broaden these results. First, including additional planetary boundaries, using consumption-based footprints for materials, land, nitrogen and phosphorus, would show whether the patterns found for climate hold across the Earth system. Second, a broader social foundation, using new data sources for education, nutrition and poverty, would reduce the dependence on four indicators. Third, case studies of the frontier countries and of countries close to the foundation could identify the policies and institutions that allow high social outcomes at low emissions. Fourth, methods that exploit policy changes or other quasi-experimental variation could move from the associations documented here towards causal estimates of how emission reductions and social outcomes interact.
